 \documentclass[12pt,twoside]{article}
\usepackage[utf8]{inputenc}
\usepackage[margin=1in,top=2.8cm,headheight=17pt]{geometry}
\usepackage{graphicx}
\usepackage{multirow}
\graphicspath{ {./images/} }
\usepackage{todonotes}
\usepackage{fancyhdr}
\pagestyle{fancy}
\fancyhead[LE, RO]{\thepage}
\fancyhead[RE]{\thechapter.\ \ \chaptername}
\fancyhead[LO]{\nouppercase{\rightmark}}
\fancyfoot{}
% za cirilicu
\usepackage[utf8]{inputenc}
% \usepackage{cmsrb} % namesti ovo kasnije da ti radi
\usepackage[OT2]{fontenc}
\usepackage[serbian]{babel}
% srpski navodnici
\def\zn{,\kern-0.09em,}
\def\zng{'\kern-0.09em'}
% forsiraj cirilicu
\newcommand{\cirilica}[1]{{\fontencoding{OT2}{\selectfont{#1}}}}
\title{Obrazovanje u Srbiji}
\author{Marina Vracharic1\and Anastasija Mandic1\and Jovana Tasovac\and Filip Manchic1}
\usepackage{titlesec}
\titleformat{\chapter}[hang]{\Huge\bfseries}{\thechapter{\ \ }}{0pt}{\Huge\bfseries}
\usepackage[hidelinks]{hyperref}
\begin{document}
\maketitle
    \renewcommand\contentsname{Sadrzhaj}
    \tableofcontents{}
    \thispagestyle{empty}
    \clearpage
    \thispagestyle{empty}
     \cleardoublepage
    

  

    \input{Onlajn nastava}
    % literatura
    \renewcommand\bibname{\cirilica{Literatura}}
    \nocite{*}
    
    
  \section{Obrazovanje kroz istoriju}
  Средњовековна српска држава је располагала одређеним школским системом. Како се уздизала и јачала све више је растао број писмених и образованих људи који су јој били потребни. У побожној православној Србији, цркве и манастири бивали су све бројнији, а посебно од времена Светог Саве, били су расадници писмености и школовања деце. Али нису били једини. На српском двору и у кућама великих властелина, деца су се учила разним вештинама: писању, читању, украшавању, ткању, трговини, занатима…све до борилачких вештина.
  \\Veliki znachaj za srpsko obrazovanje kroz istoriju nosi Resavska shkola prepisivacha. Ресавска школа постојала је током прве половине 15. века и радила је у оквиру манастира Манасије, који је подигао деспот Стефан Лазаревић, као своју задужбину. У њој су се окупљали учени монаси писци, преводиоци, књижевници, преписивачи који су украшавали рукописе и књиге, због чега је манастир представљао симбол духовности и просвећености током неколико наредних векова. Један од главних сарадника Ресавске школе био је Константин Филоsоф. По правилима ове школе радило се у Љубостињи, Хиландару, Пећкој патријаршији, Дечанима, а утицај се осећао и у Македонији, Бугарској, Румунији и Русији. \\Упркос значају и утицајима Ресавске школе, данас је нажалост у друштву раширен израз „ресавска школа“, који има потпуно другачије значење.\\ Преписи и преводи Ресавске школе сматрају се узорнима, по којима се до 18. века мерио квалитет рукописа.\\ Захваљујући раду и утицају Ресавске школе српски народ је сачувао свој језик. \\ У време турских освајања манастир Ресава делио је судбину читаве српске Деспотовине. Књижевна активност манастира ометена је најездом Турака и пропашћу деспотовине. Један део калуђера повукао се са последњим Бранковићима у Срем и наставио је ту акцију у фрушкогорским манастирима (нпр. у Крушедолу), други је отишао у Влашку, а затим у Русију, поневши са собом правописна правила и списе Константина Философа, док се најмањи део окупио у Јужној Србији, у Пећкој патријаршије и Дечанима.\\
  \\ Падом Србије под турско ропство, све се изменило. За генерацију или две нестало је српске властеле, читавог расадника писмености и учености. Оно што је код Срба било прогресивно, Турци су свесно и систематски гушили и уништавали. \\После пада српске средњовековне државе српски народ се није развијао на једном заједничком простору и под истим политичким, верским и културним утицајима.\\
  \\У 18. и 19. веку већина српског народа живи под турском и аустроугарском влашћу. Развитак школства код Срба у Хабзбуршкој монархији у 18. веку одвијао се у неколико карактеристичних етапа:\\\\
  \begin{enumerate}
      \item Sрпско-словенска етапа од 1690-1726. године;\\ Школе вероисповедног обележја почињу се оснивати одмах након доласка Срба у Аустрију. Школе су носиле назив "тривијалне школе" јер су биле организоване по узору на сличне грчке школе у Цариградској патријаршији
      \item Rуско словенска етапа од 1726-1749. године
      \item Vреме митрополита Павла Ненадовића 1749-1768 године
      \item Pериод просвећености који је био праћен терзијанским и јозефинским реформама.
  \end{enumerate}
  \\\\
  У Београду је између 1718. и 1739. постојала Мала српска школа. Почетком 18. века, у порти цркве Светог Ђорђа, 1703. год. основана је прва школа у Петроварадинском Шанцу ( Новом Саду) - Српска православна основна школа. Аврам Мразовић је 1. маја 1778. у Сомбору основао српску основну школу „Норму“, која је најстарија установа за образовање учитеља за словенско становништво на југу Европе. Карловачки митрополит Стеван Стратимировић је 1791. у Сремским Карловцима основао Карловачку гимназију. Током Првог српског устанка, у Београду је 1808. основана Велика школа заслугом Ивана Југовића и Доситеја Обрадовића. У Крагујевцу је октобра 1838. основан Лицеум Књажества сербског, прва виша школа у Србији. Он је 1841. премештен у Београд. Лицеј је постојао до 1863. када прераста у Велику школу, која је имала три факултета: Филозофски, Технички и Правни.\\\\Београдски универзитет је основан 1905. После Другог све{т}sког рата из Београдског универзитета су се издвојили Универзитет у Новом Саду (1960), Нишу (1965), Приштини (1970), Подгорици (1974) и Крагујевцу (1976).\\До примене Болоњске декларације, 2006, и свеобухватних измена образовног система, Србија је примењивала систем наслеђен из социјалистичке Југославије. Предшколско образовање било је изборно, а основно и средње образовање идентични садашњим. Применом Болоњске декларације школске 2005/06. године, бивша Диплома високог образовања је изједначена са мастером, а магистратура са прве две године докторских студија (још једна година до доктората), обоје услед једнаких година учења. Постдипломско образовање је обухватало оно што данас обухватају други и трећи део високог — магистратуру и докторат.\\
  Са образовањем деце се почиње у предшколским установама које се одржавају у локалном вртићу, те је оно први део образовања. Оно је обавезно од школске 2006/07. године. Траје бар 4 сата дневно бар 6 месеци у години уписа у први разред, за децу од пет или шест година, а циљ му је упознавање ученика са образовним системом и припрема за основну школу.\\
  Деца се уписују у основну школу са шест или седам година. Као и предшколско, и основно образовање је обавезно. Основна школа траје осам година и подељена је на два периода:
\begin{enumerate}
    \item први циклус основног образовања (од 1. до 4. разреда)
    \item други циклус основног образовања (од 5. до 8. разреда)
\end{enumerate}
\section{Onlajn nastava}
Posledica korona virusa 2019. godine jeste potpuni prelazak na onlajn nastavu shirom sveta.\\
Kao i nastava u uchionici, i onlajn nastava je zasnovana na teorijama uchenja. Polazna tachka su obrazovni ciljevi, zatim odluke o tome kako c1e znanje biti vrednovano (test, diskusija, projekat, prezentacija…). Nakon toga se biraju pedagoshki pristupi (na primer uchenje zasnovano na projektu, na timskom radu, na problemu) i alati koji c1e omoguc1iti da se ta nastava realizuje. Povratna informacija od nastavnika je kljuchan deo procesa jer bez nje uchenik ne mozhe napredovati.
 Takav rezhim nastave donosi sa sobom nekoliko dobrih i loshih strana. \\
Mnogim uchenicima bi najvec1a mana kod onlajn nastave bila to shto nisu u direktnoj interakciji sa svojim drugovima iz odeljenja i profesorima. To mozhe uzrokovati nedostatak diskusije potrebne za razumevanje lekcije. Uchenici su sada na chasu preko svojih virtuelnih mashina, shto automatski znachi da moraju imati adekvatnu opremu za prac1enje iste. \\
Velika je promena shto oni prilikom onlajn nastave ne prate predavanje u shkoli vec1 od kuc1e. Na to se mozhe gledati i sa pozitivne i sa negativne strane. Pozitivno je svakako to shto se nalaze u toplini svog doma i shto mogu sebi omoguc1iti najudobniju atmosferu za prac1enje nastave. Dok sa druge strane, nekima to ne odgovara jer taj komfor u njima budi preveliki osec1aj slobode shto prouzrokuje lenjost, kao i usamljenost.
\begin{center}
\includegraphics[width=11cm, height=8cm]{online-kurs-jezika.jpg}
\end{center}
\section{Kvalitet obrazovanja u Srbiji}
O kvalitetu obrazovanja se sve manje zakljuchuje na osnovu utisaka, lichnih stavova i uverenja, a sve vishe na osnovu objektivnih pokazatelja do kojih se dolazi istrazhivanjima i analizom rezultata istrazhivanja. Najkorisniji i najpouzdaniji su nalazi koji nisu parcijalni i izolovani, a do takvih se dolazi istrazhivanjima koja se sprovode na nacionalnom nivou i međunarodnim istrazhivanjima.

Vazhnost dobijanja pouzdanih i valjanih podataka o kvalitetu obrazovanja je jedna strana medalje. Druga je kako se dobijeni podaci i informachije koriste. Nachin na koji se to chini, publicitet koji dobijaju, karakter i nivo struchne i shire drushtvene rasprave o njima, u najvec1oj meri govore o tome kako se jedno drushtvo i drzhava odnose prema obrazovanju i koliki mu znachaj pridaju.

Istrazhivanja i poređenja kvaliteta obrazovanja u razlichitim zemljama (međunarodna istrazhivanja) su, do relativno skoro, bila veoma retka a, kako je izgledalo, praktichno skoro i nemoguc1a. Glavni razlog tome shto među zemljama postoje velike razlike u tome kako je organizovan i kako funkcionishe obrazovni sistem. Tradicija u svakoj zemlji na razlichit nachin oblikuje neke od najbitnijih karakteristika obrazovanja i obrazovnog sistema.

Ukupno gledano, rezultati nashih uchenika na TIMSS-u* su ohrabrujuc1i kada uzorak chine uchenici chetvrtog razreda, jer su njihova postignuc1a iznad međunarodnog proseka za taj uzrast, dok bi morali znachajno da se poboljshaju kad je rech o uchenicima osmog razreda. Vidljivo opadanje nivoa postignuc1a nashih uchenika sa uzrastom je veoma vazhna i upozoravajuc1a informacija za sve koji uchestvuju u obrazovnom procesu i upravljaju obrazovnim sistemom.


Vishe od trec1ine srednjoshkolaca funkcionalno je nepismeno, a istovremeno sve je vishe odlikasha i „vukovaca“ – uchenika koji imaju sve same petice.

Jaz u kome se nalazi srpsko obrazovanje mozhda se najbolje otkriva u ova dva pokazatelja.

Trenutno je u skupshtinskoj proceduri set zakona o obrazovanju, ali struchnjaci smatraju da se ni postojec1i propisi ne poshtuju dovoljno.
Kljuchni problem je to shto je status obrazovanja ,,ispod onog koji zasluzhuje''. Odgovornost prosvetara je ogromna, jer obrazuju chitave generacije, a nalaze se u podređenom polozhaju.
Nije samo novac u pitanju, ali je to definitivno jasan pokazatelj. Biti nastavnik nije vishe pozheljno zanimanje kao pre. Polovina nastavnika u osnovnim i srednjim shkolama nema stalni radni odnos, shto je ozbiljan problem.
Problem je dodatno otezhala zabrana zaposhljavanja u javnom sektoru, koja je trajala sedam godina, sve do pochetka 2021.
Pored toga, i dalje mnogi od nastavnika rade sa nepunim fondom chasova. 
\\
**TIMSS - jedan od projekata Međunarodnog udruzhenja za evaluaciju obrazovnih postignuc1a (International Association for the Evaluation of Education Achievement – IEA)
\section{Reference}
\begin{enumerate}
\item \url{https://sr.wikipedia.org/sr-el/%D0%9E%D0%B1%D1%80%D0%B0%D0%B7%D0%BE%D0%B2%D0%B0%D1%9A%D0%B5_%D1%83_%D0%A1%D1%80%D0%B1%D0%B8%D1%98%D0%B8}
\item \url{https://www.bbc.com/serbian/lat/srbija-59409115}
\item \url{https://scindeks-clanci.ceon.rs/data/pdf/0353-7129/2008/0353-71290803173S.pdf}
\item \url{https://okc.rs/sta-jeste-a-sta-nije-onlajn-nastava/}
\item \url{https://www.stat.gov.rs/sr-latn/oblasti/obrazovanje/}
\item \url{https://www.ozonpress.net/drustvo/danas-je-medjunarodni-dan-pismenosti-statistika-u-srbiji-porazavajuca/}
\end{enumerate}
\end{document}


